\subsubsection{বাংলা পেজে ভাষা ও এনকোডিং আলাদা কেন?}
\texttt{lang="bn"} জানায় নথির প্রধান ভাষা বাংলা; সহায়ক পাঠযন্ত্রের উচ্চারণ ও ভাষাভিত্তিক কাজের জন্য এটি দরকার। \texttt{charset="UTF-8"} জানায় ফাইলের বাইট কীভাবে অক্ষর হিসেবে পড়তে হবে। ভাষা ঘোষণা নিজে থেকে ফাইলের এনকোডিং বদলায় না বা লেখা অনুবাদ করে না।

\par\noindent\begin{minipage}[t]{.54\linewidth}\vspace{0pt}\small
\textbf{\color{chFourTeal}কোড: বাংলা নোটিশ}\par
\begin{chFourBanglaCode}
\begin{verbatim}
<!DOCTYPE html>
<html lang="bn">
<head>
  <meta charset="UTF-8">
  <title>বিজ্ঞান ক্লাব</title>
</head>
<body><p>সভা শুক্রবার।</p></body>
</html>
\end{verbatim}
\end{chFourBanglaCode}
\end{minipage}\hfill
\begin{minipage}[t]{.42\linewidth}\vspace{0pt}\small
\textbf{\color{chFourTeal}ফল ও করণীয়}\par
ব্রাউজারের ট্যাবে বিজ্ঞান ক্লাব; পেজে সভা শুক্রবার।\par
ফাইলটি সত্যিই UTF-8-এ সংরক্ষণ করবে। বাংলা ভেঙে এলে প্রথমে সংরক্ষণের এনকোডিং ও ঘোষণার মিল দেখবে; পরে উপযুক্ত বাংলা ফন্ট আছে কি না দেখবে।
\end{minipage}\par
শুধু \texttt{lang="bn"} দিলে ভুল এনকোডিংয়ের বাংলা ঠিক হয় না। ভাষার পরিচয় এবং বাইট পড়ে অক্ষর শনাক্ত করার নিয়ম আলাদা তথ্য।
